
\subsubsection{Key Findings}

The methodology validation experiments support the hypothesis that perplexity filtering could amplify benchmark contamination:

\begin{enumerate}
\item \textbf{CCR is a calibrated metric}: Linear scaling with injection rate ($R^2$ = 0.9998) validates CCR as a reliable measure.
\end{enumerate}

\begin{enumerate}
\item \textbf{CCR differences are detectable}: The bootstrap methodology correctly identifies statistically significant CCR differences between filtering strategies.
\end{enumerate}

\begin{enumerate}
\item \textbf{High-CCR examples show disproportionate influence}: Removal causes 1.$97\times$ greater degradation than random removal, indicating causal importance.
\end{enumerate}

\begin{enumerate}
\item \textbf{Contaminated examples are structurally irreplaceable}: 4.$1\times$ higher IFR suggests contaminated examples provide unique signal.
\end{enumerate}

\subsubsection{Limitations}

\#\#\#\# Methodology Validation Mode

All experiments were conducted using simulated contamination due to GPU infrastructure constraints. Specifically:

\begin{itemize}
\item H-E1: Validated detection mechanism; training skipped
\item H-M1: Used simulated contamination injection rather than natural perplexity-filtering differences
\item H-M2: Logic validation with synthetic accuracy patterns
\item H-M3: Eval-only mode with simulated strategy effects
\item H-C1: Simulated embeddings and TRAK scores
\end{itemize}

Full empirical validation requires GPU cluster access to train models on filtered corpora and compute real attribution scores.

\#\#\#\# Single Model Scale

Experiments are designed for 1B parameter scale (Pythia-1B). Contamination dynamics may differ at larger scales.

\#\#\#\# Proxy Perplexity Signals

Methodology validation used word-count-based proxy perplexity. Real validation requires RedPajama-V2 ccnet\_perplexity signals.

\#\#\#\# MMLU-Redux Assumption

The Amplification Index computation assumes MMLU-Redux has minimal overlap with training corpora. This assumption has not been verified via MinHash or embedding-based audit.

\subsubsection{Weaker-Than-Expected IFR-Redundancy Correlation}

The correlation between IFR and k-NN redundancy ($\rho$ = -0.11) was substantially weaker than the hypothesized threshold ($\rho$ < -0.5). Possible explanations include:

\begin{enumerate}
\item k-NN redundancy in embedding space may not capture task-relevant structure
\item Contamination may operate at sub-document (span) level rather than document level
\item Synthetic data artifacts may not reflect real contamination structure
\end{enumerate}

This finding suggests the mechanism by which contaminated examples become structurally necessary requires further investigation.

\subsubsection{Implications}

If the curation-contamination amplification effect holds under full empirical validation:

\begin{enumerate}
\item Benchmark scores may not be comparable across curation strategies
\item Decontamination has quantifiable costs (approximately $2\times$ expected degradation)
\item CCR, AI, and IFR provide tools to audit contamination effects in curation pipelines
\end{enumerate}
